\chapter{The Indological Perspective}

\section{Indology as Method: The Orientalist Background}

The first topic in the syllabus is the \textbf{Indological Perspective}, with \textbf{G. S. Ghurye} named in brackets. Indology is a \textbf{method} to study India. (A recommended recent book is \textbf{Abhijit Kundu's ``Sociology of India''}, Sage, 2021.)

It all starts with \textbf{Orientalism} --- the representation of the non-West by the Western. The British wanted to understand ``Indian culture,'' a problematic phrase since India is diverse and has no single culture.

\begin{CriticalPoint}
``Indian culture'' is a misnomer: India is not a homogeneous country with one culture; every community has its own culture.
\end{CriticalPoint}

\textbf{William Jones} and \textbf{Wilkins} are regarded as the founders of Indology. They tried to understand Indian culture through the ancient Brahmanical texts --- the Vedas, Upanishads, Bhagavad Gita, Hitopadesha, Manu Smriti.

\begin{KeyConcept}
\textbf{Hermeneutics} = the interpretation of culture through text, including the attempt to study the \textbf{intention of the author}.
\end{KeyConcept}

Jones found similarities between Sanskrit and European languages and concluded that European and Indian cultures are the same. \textbf{Max M\"uller} wrote \textbf{Sacred Books of the East}, was a professor of \textbf{philology} (comparison of languages), gave the \textbf{Aryan Invasion Theory} in a systematic form, and coined \textbf{``Aryan brotherhood.''}

\begin{KeyConcept}
\textbf{Phase of Indomania} --- the three key figures are \textbf{William Jones, Wilkins and Max M\"uller}: all in love with Indian culture, all finding similarities between Indian and European culture, and finally calling the link a ``brotherhood.''
\end{KeyConcept}

\begin{QuickRevision}
\begin{itemize}
\item Indology = a method; first syllabus topic (G. S. Ghurye).
\item Founders: Jones \& Wilkins; hermeneutics; Max M\"uller's Aryan Invasion Theory \& Aryan brotherhood.
\item Indomania = love of Indian culture (Jones, Wilkins, Max M\"uller).
\end{itemize}
\end{QuickRevision}

\section{From Indomania to Indophobia}

\textbf{Indophobia} is the phase of criticising Indian culture to justify colonial rule. Its notable figures are \textbf{Max Weber} and \textbf{Karl Marx} --- and, later, Max M\"uller himself.

\begin{itemize}
\item \textbf{Max Weber} is an Orientalist: he represented Hinduism in a specific manner, and his account fed the colonial justification for rule.
\item \textbf{Karl Marx} criticised Indian culture: the Indian village was a \textbf{static unit} that does not change, hence no progress. (India has no class structure, only caste; religion prevents antagonism --- Dalits respect Brahmins --- so society remains static; whatever change came to India came from the British: modern education, middle class, English education, the telegram and the postal network.)
\item \textbf{Max M\"uller} later turned to a \textbf{theory of racial division} --- the indigenous race is inferior and must be civilised by the Aryan race.
\end{itemize}

\begin{KeyConcept}
\textbf{Rule through knowledge:} colonial rule was justified through knowledge, not just military force. (Weber's power and authority always require people's acceptance for legitimacy.)
\end{KeyConcept}

\begin{CriticalPoint}
The shift from Indomania to Indophobia is called the \textbf{European U-turn}, developing gradually with the consolidation of British rule.
\end{CriticalPoint}

\begin{QuickRevision}
\begin{itemize}
\item Indophobia = criticising Indian culture to justify rule (Weber, Marx, later Max M\"uller).
\item Rule through knowledge; the European U-turn.
\end{itemize}
\end{QuickRevision}

\section{Nationalist / Indigenous Indology (Bombay School)}

In reaction to Indophobia, \textbf{nationalist historiography} developed, to restore the glory of India --- and with it \textbf{cultural nationalism}.

Two institutions: the \textbf{Asiatic Society} (William Jones --- to study Hindu/Sanskritic culture) and the \textbf{Bhandarkar Research Institution} (Bombay --- to criticise the Orientalist approach and restore India's past glory).

\begin{ExampleBox}
The quarrel is about ``\textbf{who is original}'': the British claimed the Hindus are like them (a duplicate); nationalist Indology claims the \textbf{Vedic Aryans are original} and the Europeans the duplicate. \textbf{Dayanand Saraswati} put it bluntly: the \textbf{Vedas are the true source of knowledge}, not the Bible.
\end{ExampleBox}

Nationalist-Indology figures include \textbf{Ranade, Bal Gangadhar Tilak, Savarkar, G. S. Ghurye and Iravati Karve} --- many of them extreme Hindus whose symbols (the Ganga, etc.) later became \textbf{Hindutva}, which discouraged Muslims from joining the nationalist movement (nationalism was equated with Hinduism).

\begin{KeyConcept}
\textbf{Karel Upadhyay} (an Indian-American sociologist) says Ghurye's \textbf{sociological imagination} can be traced to \textbf{three factors: British Orientalism, Diffusionism, and Nationalism.}
\end{KeyConcept}

\begin{KeyConcept}
\textbf{Brahminization} is Ghurye's term (later popularised by M. N. Srinivas as \textbf{Sanskritization}): the process of Aryanising/Hinduising the South and the tribals.
\end{KeyConcept}

\subsection{The Bombay school}

\begin{itemize}
\item The indigenous/nationalist Indology is also called the \textbf{Bombay school} (like the Frankfurt school or the Chicago school). It developed in reaction to British Orientalism, Nationalism and Indophobia.
\item Led by \textbf{G. S. Ghurye} (founder), with \textbf{Iravati Karve, Bankim Chattopadhyay} and \textbf{B. D. Shankar}.
\item It tried to \textbf{restore and assert India's true place in human history}, and to \textbf{neutralise the Western Indologists (Indophobics)} who were portraying a wrong picture of Indian society.
\item It claims the \textbf{superiority of Vedic Aryans over European Aryans} --- everything in Europe had already been invented in the Indian past, and the Vedas are the true source of knowledge.
\end{itemize}

This social thought led to \textbf{cultural nationalism} --- seen in the \textbf{Dharma Sabha, Brahmo Samaj and Arya Samaj}, all of which trace Indian roots back to Vedic Hindu Aryan culture (to the Gita and the Vedas).

\begin{CriticalPoint}
In this discourse Ghurye saw \textbf{India as the product of the Vedic period} and \textbf{ignored the contributions of Islamic and British rulers}. Equivalences follow: Indian culture = Vedic culture; Indian philosophy = Vedanta; Indian religion = Hinduism; Aryans = the higher castes (Brahmins). The result --- ``\textbf{India is Hindu and Hindu is India.}''
\end{CriticalPoint}

\begin{QuickRevision}
\begin{itemize}
\item Nationalist Indology = Bombay school (reaction to Indophobia); Bhandarkar Research Institution.
\item Ghurye (founder) with Iravati Karve, Bankim Chattopadhyay, B. D. Shankar.
\item Upadhyay's three factors: Orientalism, Diffusionism, Nationalism; Brahminization, later Sanskritization.
\end{itemize}
\end{QuickRevision}

\section{G. S. Ghurye: Cultural Unity \& Composite Culture}

Ghurye is one of the most important pillars: M. N. Srinivas, Andre Beteille and A. R. Desai were all, in one way or another, influenced by him.

\textbf{D. P. Mukerji} held that ``you cannot be an Indian sociologist unless you know Sanskrit,'' and regarded Ghurye as the Indian sociologist who knows Sanskrit. Ghurye established the \textbf{Department of Indology at Bombay University (1919)}.

\begin{KeyConcept}
Ghurye's core dilemma: should Indian sociology follow Western frameworks, or be indigenous? His answer --- you cannot have your own sociology without understanding \textbf{Hindu culture} (the culture of the texts).
\end{KeyConcept}

Like the Europeans, Ghurye compared civilisations (Indian, Egyptian, European): joint family in India/Egypt vs nuclear family in Europe; jointness vs individualism. ``You cannot have sociology without comparison --- comparative method is sociology itself.''

\begin{KeyConcept}
Ghurye believed \textbf{Indian culture is one (homogeneous)}, and that this \textbf{cultural unity} can be traced to the \textbf{ancient texts}. Indian culture = \textbf{Hindu Vedic Aryan culture}, formed in the Indo-Gangetic plains and diffused to the South.
\end{KeyConcept}

\begin{CriticalPoint}
To privilege Sanskrit (and the Vedas) is to privilege one culture over another: India had the \textbf{Sangam literature} too, and Veda is not Sangam, Sangam is not Veda. Equating the Sanskrit/Vedic culture with Indian culture as such is the core Indological problem.
\end{CriticalPoint}

Institutions that created cultural unity are part of Indian culture. The transmission from north to south happened through:
\begin{itemize}
\item \textbf{Brahmi script + Sanskrit language} --- Sanskrit was a medium of conversation and inter-regional contact; many languages are Brahmi-scripted and Sanskrit-derived (e.g.\ Marathi).
\item \textbf{Marital alliances} --- a political instrument through which royal families (North-West/Sindh, Kamarupa/Assam, Vidarbha, Dravid desh) came into contact, spreading culture.
\item \textbf{Ashwamedha sacrifice} --- the horse, a symbol of monarchic power, wandered all over India and provided \textbf{ritual unity}.
\item \textbf{Indian sadhus and saints} --- travelling sages transmitting Vedic culture (Ghurye's books \textbf{Gods and Men} and \textbf{Indian Sadhus}). Ghurye names \textbf{Sage Agastya, Shankaracharya and God Skanda/Kandha} as carriers of Vedic Aryanism to the South.
\end{itemize}

\begin{QuickRevision}
\begin{itemize}
\item Ghurye: cultural unity traceable to ancient texts; composite culture.
\item D. P. Mukerji on Sanskrit; Bombay University 1919.
\item Diffusion mechanisms: Brahmi + Sanskrit, marital alliances, Ashwamedha, sadhus/saints.
\item Sangam is not Veda --- privileging Sanskrit privileges one culture over another.
\end{itemize}
\end{QuickRevision}

\section{The Triad}

\begin{KeyConcept}
The \textbf{Triad} is the idea of ``three'' in everything in Hinduism --- a \textbf{set of values} that forms the composite culture of the Aryan Vedic Hindus.
\end{KeyConcept}

Its components:
\begin{itemize}
\item \textbf{Three gods} --- Brahma, Vishnu, Mahesh.
\item \textbf{Purusharth} (purposes of life) --- \textbf{Dharma} (duty), \textbf{Artha} (economic condition), \textbf{Kama} (passion), \textbf{Moksha} (salvation).
\item \textbf{Ashrams} (stages of life) --- \textbf{Brahmacharya} (student), \textbf{Grihastha} (householder), \textbf{Vanaprastha} (anchorite), \textbf{Sanyasa} (renunciate); Ghurye clubs the last two to make a triad.
\item \textbf{Trai Vidya} (threefold knowledge) --- \textbf{Rig Veda} (rituals/code of conduct), \textbf{Yajur Veda} (war and the warrior class), \textbf{Sam Veda} (art, music, dance); the Atharva Veda is excluded.
\item \textbf{Three qualities} --- Sattva, Rajas, Tamas.
\item \textbf{Threefold ethics} --- \textbf{Dam} (self-control), \textbf{Daan} (charity), \textbf{Daya} (compassion).
\item \textbf{Three religious duties} --- \textbf{Yajna} (sacrifice), \textbf{Adhyayan} (study), \textbf{Dana} (charity).
\item \textbf{Three Varnas} --- Brahmin, Kshatriya, Vaishya (written ``Vis'' in the original Rig Veda; the Shudra is not mentioned in the original Veda).
\end{itemize}

\begin{ExampleBox}
The roots of corporate social responsibility (CSR) can be traced to the \textbf{Daan} philosophy, and Gandhi's idea of \textbf{trusteeship} comes from Daan --- unique to Indian culture.
\end{ExampleBox}

\begin{DataFact}
The ages of the ashrams vary by Varna: Brahmacharya begins at 5--6 for Brahmins, 7--8 for Kshatriyas, 9--10 for Vaishyas, and is \textbf{absent for the Shudra}.
\end{DataFact}

\begin{QuickRevision}
\begin{itemize}
\item Triad = the ``three'' in everything (set of values that forms composite culture).
\item Components: 3 gods, Purusharth, ashrams, Trai Vidya, 3 qualities, 3 ethics, 3 duties, 3 Varnas.
\item CSR/Gandhi's trusteeship come from the Daan philosophy.
\end{itemize}
\end{QuickRevision}

\section{Diffusionism \& the Aryanization of the South}

To explain how one culture spread, Ghurye used \textbf{Diffusionism} --- a theory associated with his PhD supervisor \textbf{W. H. R. Rivers}.

\begin{center}
\begin{tabularx}{\textwidth}{YYY}
\rowcolor{AccentDark}
\thead{Attribute} & \thead{Diffusionism} & \thead{Evolutionism} \\
Origin & One --- culture begins at a single point (Egypt/Persia) and diffuses outward, taking different forms & Many --- each culture has its own starting point \\
Direction & Culture travels out from one origin & Cultures converge on one end-point (the scientific age) \\
Core claim & The origin is one & The development/end is one \\
\end{tabularx}
\end{center}

The Aryan culture travelled from Egypt/Persia to the Indo-Gangetic plains (and separately to Europe), where it became \textbf{Vedic Hinduism}. \textbf{Caste is described as the original creation of the Indo-Gangetic plains}, from where it diffused to the South.

\begin{KeyConcept}
\textbf{Aryanization of the South} --- Ghurye says \textbf{Agastya, Shankaracharya and Skanda (Kandha)} carried Vedic Aryanism to the South. The South is thus not ``different,'' only an Aryanised copy.
\end{KeyConcept}

\begin{itemize}
\item \textbf{Skanda = Kartik = Subramaniam (Murugan)} in the South --- the same god in a diffused, Dravidian form.
\item \textbf{Krishna} in the North is fair/blue, full-bodied, playing the flute; in the South he is black, handless, without flute. \textbf{Jagannath} is another form (avatar) of Vishnu.
\item \textbf{Buddha and Mahavir are claimed as avatars of Vishnu} --- the Hinduization of non-Hindus.
\end{itemize}

\begin{CriticalPoint}
The Aryan timeline is itself open to the \textbf{sociology of knowledge (Karl Mannheim)}: the carbon dating is ``traced'' to about 1500 BC, but one already influenced by the Aryan Invasion theory may look only that far --- prior theory shapes research.
\end{CriticalPoint}

\subsection{Syncretism}

\textbf{Syncretism} = continuous interaction between North and South, and between folk and elites --- e.g.\ \textbf{Kabir} seeking similarities between Hinduism and Islam to show we are culturally one.

\begin{ExampleBox}
The assimilation of tribal gods into Hinduism: \textbf{Hanuman} (monkey form) and \textbf{Ganesh} (elephant form) show that animal worship (animism/animatism) was absorbed into Hindu culture.
\end{ExampleBox}

Ghurye extends this: Jainism (Mahavir), Buddhism (Buddha) and Sikhism (Guru Nanak --- who travelled as far as Multan) are all, in his reading, \textbf{part of Hindu culture in a diffused form}. J. P. S. Oberoi calls Buddhism, Jainism and Sikhism the \textbf{``Hindu protestant religions.''}

\textbf{Louis Dumont} is also an Indologist (using structuralism as a method): all over India he found one common thing --- the \textbf{Brahmin--untouchable binary of purity and impurity} --- another expression of cultural unity.

\begin{QuickRevision}
\begin{itemize}
\item Diffusionism (origin one) vs Evolutionism (end one); theory of W. H. R. Rivers.
\item Caste = original creation of the Indo-Gangetic plains; diffused to the South.
\item Aryanization of the South (Agastya, Shankaracharya, Skanda); syncretism; Louis Dumont's purity--impurity.
\item Karl Mannheim's sociology of knowledge: prior theory shapes research (carbon dating).
\end{itemize}
\end{QuickRevision}

\section{Ghurye on Caste}

\begin{KeyConcept}
Ghurye: ``\textbf{Caste is a Brahmanical child of Indo-Aryan culture, cradled in the land of Ganga and Jamuna, and then transferred to other parts of the country through diffusion.}'' (A line to reproduce in the exam.)
\end{KeyConcept}

\begin{CriticalPoint}
Because Ghurye accepted the Aryan Invasion (a racial) theory, he subscribes to the \textbf{racial theory of caste} --- caste is the child of race. Hence his first book on caste was \textbf{``Caste and Race.''}
\end{CriticalPoint}

\begin{KeyConcept}
``\textbf{Race produced caste, but occupation stabilised it.}'' So Ghurye combines the \textbf{racial} and \textbf{occupational} theories of origin.
\end{KeyConcept}

\begin{DataFact}
Ghurye's book on caste went through \textbf{5 editions} (with changing titles): Caste and Race, then Caste and Class, then Caste, Class and Occupation, and the last edition (1969) on what is happening to caste today. The first edition was \textbf{1932}.
\end{DataFact}

\begin{QuickRevision}
\begin{itemize}
\item Caste = Brahmanical child of Indo-Aryan culture, cradled in the Ganga--Jamuna land, transferred by diffusion.
\item Race produced caste, occupation stabilised it (racial + occupational).
\item First book ``Caste and Race'' (1932; 5 editions).
\end{itemize}
\end{QuickRevision}

\section{Joint Family, Gotra \& Charana}

\subsection{Joint family}

\begin{KeyConcept}
\textbf{Joint family:} four or more people living under one roof, sharing food and property --- the \textbf{Hindu Undivided Family (HUF)}.
\end{KeyConcept}

Ghurye traces the joint family to the epics (the family-centric Ramayana/Mahabharata, five brothers together), contrasting it with today's hero-centric, individualistic serials. The British codified Hindu family law on the basis of these ancient Sanskrit texts --- laws that persist in the HUF.

\begin{CriticalPoint}
The fear among minorities is that the \textbf{Uniform Civil Code} is the modern, secular avatar of the \textbf{Aryanisation/Hinduisation of the non-Hindu}.
\end{CriticalPoint}

\subsection{Gotra \& Charana}

Ghurye wrote \textbf{``Brahminical Institutions: Gotra and Charana.''}

\begin{KeyConcept}
\begin{itemize}
\item \textbf{Gotra} derives from caste; there are \textbf{seven gotras} (from the seven/sapta rishis), whose \textbf{36 children} form the \textbf{pravara}. In caste we practise \textbf{endogamy}; in gotra we practise \textbf{exogamy} (sapind exogamy) --- one cannot marry within the same gotra.
\item \textbf{Charana} is the \textbf{guru--shishya relationship}, taught in the \textbf{gurukul} (the school revolving around the guru), a \textbf{dyadic} bond --- unlike today, where a teacher is honoured only once a year (Teacher's Day) and coaching batches run into the hundreds.
\end{itemize}
\end{KeyConcept}

\begin{KeyConcept}
\textbf{Integration} is the common idea running through Ghurye's work: \textbf{caste divides, gotra unites}; the joint family integrates through jointness. Even the caste system, in his reading, is an \textbf{integrative system} --- a Brahmin and a Dalit both have caste and are thus integrated.
\end{KeyConcept}

\begin{QuickRevision}
\begin{itemize}
\item Joint family = HUF (4+ under one roof).
\item Gotra: endogamy in caste, exogamy in gotra; Charana = guru--shishya (gurukul).
\item Integration is the running theme: caste divides, gotra unites.
\end{itemize}
\end{QuickRevision}

\section{The Indian Knowledge System (Vidyas)}

\begin{KeyConcept}
The \textbf{Vidyas} are the Indian system of knowledge --- a compendium of \textbf{32 texts/topics} (ritual, martial, practical sciences and philosophy: Mimamsa, Samkhya, Vaisesika; and the \textbf{18 Puranas}).
\end{KeyConcept}

\begin{CriticalPoint}
The Vidyas were \textbf{limited to elites} --- the Kauravas and Pandavas were elites; the masses were excluded. Inequality is visible even in the Mahabharata.
\end{CriticalPoint}

\begin{QuickRevision}
\begin{itemize}
\item Vidyas = compendium of 32 texts/topics + 18 Puranas.
\item Limited to elites --- inequality even in the Mahabharata.
\end{itemize}
\end{QuickRevision}

\section{Subject Matter: Social History}

Ghurye was one of the first to establish \textbf{social anthropology} in India. Asked what Indian sociology should study, he answered: \textbf{social history} --- history on social lines, where facts are integrated into a theoretical framework.

\begin{DataFact}
In \textbf{1931} Ghurye gave a lecture on the role of history and sociology, arguing that \textbf{distinguishing history from sociology would be fruitless} in India --- sociology cannot be visualised without historical (ancient Sanskritic) roots.
\end{DataFact}

\begin{KeyConcept}
Ghurye's three subject matters of sociology: (1) \textbf{social history}, (2) \textbf{comparison of cultures} (comparative method = sociology), (3) \textbf{history of civilisations}.
\end{KeyConcept}

\begin{CriticalPoint}
Ghurye is sometimes called the \textbf{``Indian Durkheim,'' } but this is imprecise: Durkheim studied \textbf{economic} integration (mechanical/organic solidarity), while Ghurye studied integration arising from \textbf{culture}.
\end{CriticalPoint}

His students carried the project forward: \textbf{M. N. Srinivas} studied the \textbf{Coorg (Kurg)} people --- on Ghurye's advice, because Ghurye believed the Coorgs were descendants of the Egyptians; \textbf{Iravati Karve} compared kinship/marriage patterns across North, Central and South India to find the underlying cultural unity.

\begin{KeyConcept}
\textbf{Social change in India}, for Ghurye, has been a constant process of \textbf{interaction between Aryans and non-Aryans}, and the gradual \textbf{assimilation} of non-Aryans into the Vedic Hindu Aryan fold.
\end{KeyConcept}

\begin{QuickRevision}
\begin{itemize}
\item Subject matter = social history; 1931 lecture (history--sociology split is fruitless).
\item Three subject matters: social history, comparison of cultures, history of civilisations.
\item ``Indian Durkheim'' (imprecise); students: Srinivas (Coorg), Iravati Karve (kinship).
\end{itemize}
\end{QuickRevision}

\section{Assumptions \& Features of Indology}

\subsection{Assumptions}

\begin{itemize}
\item India is a \textbf{culturally united civilisation} with a \textbf{homogeneous} population.
\item Indian culture can be understood \textbf{only} by referring to the ancient Sanskrit texts.
\item These texts reveal the \textbf{underlying ideas and values beneath diversity} --- what unites the seemingly different cultures (e.g.\ Dravidian and Brahmanical).
\end{itemize}

\subsection{Features}

\begin{itemize}
\item It is a \textbf{book view} of Indian society --- knowledge of culture comes only through text.
\item Reference to ancient texts establishes \textbf{continuities between present and past} (hermeneutics).
\item The texts help formulate an \textbf{ideal type} against which divergences in the present are understood.
\end{itemize}

\begin{QuickRevision}
\begin{itemize}
\item Assumptions: India homogeneous; texts-only; texts reveal unity beneath diversity.
\item Features: book view; continuity present--past; ideal type.
\end{itemize}
\end{QuickRevision}

\section{Indological Explanations of Contemporary Issues}

The Indologist explains the present by tracing it to the ancient texts. The teacher gives four examples (for answers), each showing either convergence or divergence between text and context:

\subsection{Untouchability}

\begin{ExampleBox}
In the Ramayana, Ram ate the berries (jhootha) of \textbf{Shabri} and \textbf{Hanuman} sat beside Ram --- so untouchability did not exist then; it was \textbf{introduced later by the foreign Aryans}.
\end{ExampleBox}

\subsection{Tribals \& environmentalism}

\begin{ExampleBox}
\textbf{Nishad} and the \textbf{Vanar Sena} helped Ram cross the river --- an integrated relationship with tribals. Ram bows before the river, showing an \textbf{indigenous environmentalism} that needed no Western education.
\end{ExampleBox}

\begin{ExampleBox}
The \textbf{Earth Goddess} is the earliest form of nature worship --- worshipping nature is worshipping woman, since both are producers and both reproduce themselves. Sita is shown close to nature, and Indian environmentalism was \textbf{initiated by women} (women are more environmentalist than men in India), unlike the West where it was initiated by men. Likewise, gender disparity in \textbf{property ownership} (land held by the king, never the queen) can be traced to the ancient texts.
\end{ExampleBox}

\subsection{Women's position}

\begin{ExampleBox}
\textbf{Iravati Karve's ``Yugant''} (1968) --- the first female recipient of the Sahitya Akademi Award --- is a feminist interpretation of the Mahabharata: \textbf{Bhishma Pitamah} objectified women (Amba, Ambika, Ambalika); \textbf{Kunti} (who blindfolded herself), \textbf{Gandhari} and \textbf{Sita} (whose \textbf{agni pariksha} is a virginity test) show that women were never safe, and that the sacrificial woman is the ideal.
\end{ExampleBox}

\subsection{Communalism}

\begin{ExampleBox}
\textbf{T. N. Madan}, reading \textbf{Kalhana's Rajatarangini} (Kalhana is called the first official historian of India), traces the coexistence and tension between Hindu and Muslim empires in Kashmir to explain the roots of communalism.
\end{ExampleBox}

\begin{CriticalPoint}
\begin{itemize}
\item Indology can be combined with Western theories: Iravati Karve uses \textbf{feminism}; D. P. Mukerji fused \textbf{Marxism} with Indology. Western theories are not opposite to Indology.
\item Ghurye's Indology is distinct in three ways: he reads the \textbf{original Sanskrit} (not translations); he uses the \textbf{diffusionist method}; and he uses Indology as a \textbf{nationalist} project to neutralise Western Indologists.
\end{itemize}
\end{CriticalPoint}

\begin{QuickRevision}
\begin{itemize}
\item Examples: untouchability (Shabri), tribals (Nishad/Vanar Sena), women (Yugant; Sita's agni pariksha), communalism (Rajatarangini).
\item Indology can be used with Western theories (Iravati Karve = feminism; D. P. Mukerji = Marxism).
\end{itemize}
\end{QuickRevision}

\section{A Note on the Critique of Indology}

The teacher flagged that the \textbf{critique of Indology} is important but deferred it to the following class (``Tomorrow we will begin with Structural Functionalism''). The critique is therefore referenced but \textbf{not delivered} in these transcripts.

\begin{QuickRevision}
\begin{itemize}
\item Critique of Indology was referenced but deferred --- not delivered in these transcripts.
\end{itemize}
\end{QuickRevision}
